\section{DDS: El Middleware de Comunicaciones}

DDS (Data Distribution Service) es un estándar de middleware publicado por el Object Management Group (OMG) en 2004~\cite{omg2015dds}. Define un modelo de comunicación pub-sub completamente distribuido: no hay servidor central, no hay broker, no hay proceso coordinador. Cada participante anuncia su presencia en la red y descubre a los demás de forma autónoma.

ROS~2 adoptó DDS como su capa de transporte en lugar de construir una solución propia. La razón fue pragmática: DDS ya resolvía los problemas que ROS~1 no podía resolver (descubrimiento distribuido, calidad de servicio, tiempo real, seguridad), y era un estándar de industria con múltiples implementaciones maduras probadas en sistemas críticos como aviónica y defensa~\cite{macenski2022robot}.

\subsection{Descubrimiento distribuido}

El protocolo subyacente de DDS es RTPS (Real-Time Publish Subscribe Protocol). Cuando un nodo ROS~2 arranca, su implementación DDS inicia el proceso de descubrimiento en dos fases.

La primera fase es SPDP (Simple Participant Discovery Protocol): el participante difunde periódicamente un anuncio de su existencia por multicast UDP en la red local. Otros participantes reciben el anuncio y añaden al nuevo participante a su lista de participantes conocidos.

La segunda fase es SEDP (Simple Endpoint Discovery Protocol): una vez que dos participantes se conocen, intercambian información sobre sus endpoints (publicadores y suscriptores) mediante comunicación punto a punto. Si un publicador y un suscriptor comparten el mismo nombre de tópico y tipo de mensaje compatible, DDS establece el canal de comunicación automáticamente.

El resultado: los nodos ROS~2 se descubren y conectan entre sí sin ninguna configuración manual de IPs, puertos, ni direcciones. Basta con que estén en la misma red y compartan el mismo \texttt{ROS\_DOMAIN\_ID}.

\subsection{El ROS\_DOMAIN\_ID}

El \texttt{ROS\_DOMAIN\_ID} es un entero entre 0 y 232 que delimita el espacio de descubrimiento. Solo los participantes con el mismo Domain ID se descubren entre sí. Esto permite correr múltiples sistemas ROS~2 independientes en la misma red física sin interferencia: un robot de prueba con \texttt{ROS\_DOMAIN\_ID=0} y un robot de producción con \texttt{ROS\_DOMAIN\_ID=1} son invisibles el uno para el otro aunque compartan la misma red WiFi.

En trabajo con múltiples computadoras (por ejemplo, robot + estación base), basta con que ambas tengan el mismo \texttt{ROS\_DOMAIN\_ID} exportado y estén en la misma subred. No se requiere configuración adicional.

\subsection{Políticas QoS}

Las políticas de calidad de servicio son la razón principal por la que DDS fue elegido sobre protocolos más simples. Cada publicador y suscriptor puede configurarse independientemente. Las más relevantes en la práctica son las siguientes.

\textbf{Reliability.} Define si la entrega de mensajes está garantizada. \texttt{RELIABLE} asegura que cada mensaje llegue al suscriptor, retransmitiéndolo si es necesario. \texttt{BEST\_EFFORT} envía sin confirmación: si el mensaje se pierde en la red, se pierde. Para streams de sensores de alta frecuencia donde el dato más reciente es siempre más valioso que un dato atrasado, \texttt{BEST\_EFFORT} es la elección correcta. Para comandos de seguridad o configuración, \texttt{RELIABLE} es obligatorio.

\textbf{Durability.} Define qué ocurre con los mensajes publicados antes de que el suscriptor exista. \texttt{VOLATILE} descarta todos los mensajes anteriores: el suscriptor solo recibe lo publicado después de suscribirse. \texttt{TRANSIENT\_LOCAL} guarda los últimos N mensajes y los entrega al suscriptor cuando este aparece. Esto es útil para tópicos de configuración que se publican una sola vez al inicio del sistema.

\textbf{History.} Define cuántos mensajes pasados se conservan en cola cuando el suscriptor no los procesa suficientemente rápido. \texttt{KEEP\_LAST(N)} conserva los últimos N mensajes. \texttt{KEEP\_ALL} conserva todos, limitado por la memoria disponible.

\textbf{Deadline.} Permite especificar que un publicador debe publicar al menos cada cierto período. Si el publicador no cumple el deadline, el middleware notifica al suscriptor. Útil para detectar nodos que han dejado de funcionar.

\textbf{Liveliness.} Define cómo se declara y verifica que un publicador sigue activo. Útil para detectar fallos de nodos sin necesitar un mecanismo de heartbeat manual.

La Tabla~\ref{tab:qos} resume las políticas más usadas con sus valores típicos en sistemas robóticos.

\begin{table}[H]
  \centering
  \caption{Políticas QoS y sus valores típicos en robótica.}
  \label{tab:qos}
  \begin{tabular}{@{}lll@{}}
    \toprule
    Política & Sensor / alta frec. & Comando / config. \\
    \midrule
    Reliability  & \texttt{BEST\_EFFORT} & \texttt{RELIABLE} \\
    Durability   & \texttt{VOLATILE}     & \texttt{TRANSIENT\_LOCAL} \\
    History      & \texttt{KEEP\_LAST(1)} & \texttt{KEEP\_LAST(10)} \\
    \bottomrule
  \end{tabular}
\end{table}

\subsection{La capa RMW: abstracción de la implementación}

DDS es un estándar, no una implementación. Existen múltiples implementaciones del estándar, cada una con distintas características de rendimiento, licencia y plataformas soportadas. Para que ROS~2 no quede atado a ninguna implementación específica, existe la capa RMW (ROS Middleware), que define una interfaz C abstracta que cualquier implementación DDS puede satisfacer~\cite{macenski2022robot}.

Las implementaciones principales son las siguientes. \textbf{Fast DDS} (eProsima) es la implementación por defecto en ROS~2 Humble. Es de código abierto, tiene excelente soporte de la comunidad, y es la más documentada en el contexto de ROS~2~\cite{eprosima2023fastdds}. \textbf{Cyclone DDS} (Eclipse Foundation) tiene menor latencia en redes locales y es la alternativa más popular para sistemas donde la latencia es crítica~\cite{cyclonedds2023}. \textbf{Zenoh} es un protocolo más moderno que extiende DDS con soporte nativo para comunicación a través de internet y dispositivos con recursos limitados; está soportado oficialmente desde ROS~2 Jazzy~\cite{zenoh2023}.

Cambiar la implementación no requiere recompilar el código de la aplicación. Basta con exportar la variable de entorno antes de lanzar el sistema, como muestra la Terminal~\ref{lst:rmw_impl}:

\begin{figure}[H]
\centering
\begin{lstlisting}[style=bash, caption={Selección de implementación DDS en tiempo de ejecución.}, label={lst:rmw_impl}]
export RMW_IMPLEMENTATION=rmw_cyclonedds_cpp
ros2 launch simulador demo_pva.launch.py
\end{lstlisting}
\end{figure}

Esta flexibilidad es uno de los puntos más fuertes de la arquitectura de ROS~2: el código de aplicación escrito con \texttt{rclpy} o \texttt{rclcpp} es completamente agnóstico respecto a qué implementación DDS está corriendo por debajo.
